% !TeX root = ../master.tex

\lecture{1}{2026-08-25}{Introduction to Heat Transfer}

\begin{problemwithimage}{p1-31}{1--31}
  A hair dryer is basically a duct in which a few layers of electric resistors are placed.
  A small fan pulls the air in and forces it to flow over the resistors where it is heated.
  Air enters a \qty{900}{\W} hair dryer at \qty{100}{\kPa} and \qty{25}{\celsius}, and leaves at \qty{50}{\celsius}.

  The cross-sectional area of the hair dryer at the exit is \qty{60}{\cm\squared}.
  Neglecting the power consumed by the fan and the heat losses through the walls of the hair dryer, determine

  \begin{subproblems}
    \subproblem the volume flow rate of air at the inlet and
    \subproblem the velocity of the air at the exit.
  \end{subproblems}
\end{problemwithimage}

\begin{givens}
  \dot{W}_e = \qty{900}{\W} & Power consumption of the hair dryer \\
  P_1 = \qty{100}{\kPa}  & Pressure at the inlet \\
  T_1 = \qty{25}{\celsius}  & Temperature at the inlet \\
  T_2 = \qty{50}{\celsius}  & Temperature of the outlet \\
  A_2 = \qty{60}{\cm^2}  & Area of the outlet
\end{givens}

\begin{assumptions}
  \item The air is an ideal gas
  \item The system is at steady state, $\frac{\mathrm{d}}{\mathrm{d}t} = 0$
  \item The fan does not consume power and no heat is lost through the wall
  \item $c_p$ is assumed constant such that the enthalpy change is simply $c_p (T_2 - T_1)$
  \item The change in pressure and potential and kinetic energy of the air is negligible
\end{assumptions}

\begin{derivation}

  \derivationpart

  Using the ideal gas law
  \begin{equation*}
    p = \rho R T
  \end{equation*}
  we can solve for the density of the air at the inlet and outlet respectively, remembering we assume the pressure is the same at both places
  \begin{align*}
    \rho &= \frac{p}{R T} \\
  \rho_1 &= \frac{\qty{100}{\kPa}}{\qty{298,15}{\K} \cdot \qty{0,287}{\kPa.\m\cubed\per\kg.\K}} = \qty{1,169}{\kg\per\m^3} \\
  \rho_2 &= \frac{\qty{100}{\kPa}}{\qty{323,15}{\K} \cdot \qty{0,287}{\kPa.\m\cubed\per\kg.\K} } = \qty{1,078}{\kg\per\m^3} 
  .\end{align*}
  As the entirety of the heat supplied by the heating element is assumed to go into the air stream the energy balance becomes
  \begin{equation*}
    \dot{W}_e = \dot{m} c_p (T_2 - T_1) = \rho_1 Q_{V} c_p (T_2 - T_1)
  .\end{equation*}
  We can now solve for the volume flow rate, $Q_V$ as
  \begin{equation*}
    Q_V = \frac{\dot{W}_e}{(T_2 - T_1) c_p \cdot \rho_1} = \frac{\qty{900}{\W}}{\qty{25}{\K} \cdot \qty{1007}{\J\per\kg\per\K} \cdot \qty{1,169}{\kg\per\m\cubed} } =\qty{0.03058}{\m\cubed\per\s} 
  .\end{equation*}
 
  \begin{subresult}
    The volume flow rate through the hair dryer is
    \begin{equation*}
      Q_V = \qty{0,031}{\m\cubed\per\s} 
    .\end{equation*}
  \end{subresult}

  \derivationpart

  We can use a similar formula as above, however, this time dividing by the outlet area as well to obtain the outlet flow speed:
  \begin{equation*}
    V_2 = \frac{\dot{W}_e}{(T_2 - T_1) c_p \cdot \rho_2 \cdot A_2} = \frac{\qty{900}{\W}}{\qty{25}{\K} \cdot \qty{1007}{\J\per\kg\per\K} \cdot \qty{1,078}{\kg\per\m\cubed} \cdot \qty{60}{\cm\squared} } = \qty{5,526}{\m\per\s} 
  .\end{equation*}
  \begin{subresult}
    The outlet flow speed of the air is
    \begin{equation*}
      V_2 = \qty{5,526}{\m\per\s} 
    .\end{equation*}
  \end{subresult}
\end{derivation}

\begin{problemwithimage}{p1-59}{1--59}
  One way of measuring the thermal conductivity of a material is to sandwich an electric thermofoil heater between two identical rectangular samples of the material and to heavily insulate the four outer edges, as shown in the figure.
  Thermocouples attached to the inner and outer surfaces of the samples record the temperatures.
  
  During an experiment, two \qty{0,5}{\cm} thick samples $\qty{10}{\cm} \times \qty{10}{\cm}$ in size are used.
  When steady operation is reached, the heater is observed to draw \qty{25}{\W} of electric power, and the temperature of each sample is observed to drop from \qty{82}{\celsius} at the inner surface to \qty{74}{\celsius} at the outer surface.
  Determine the thermal conductivity of the material at the average temperature.
\end{problemwithimage}

\begin{givens}
  h = \qty{0,5}{\cm} & Height of specimen \\
  w = \qty{10}{\cm} & Width of specimen \\
  \dot{Q}_{\text{heater}} = \qty{25}{\W} & Power-consumption of heater \\
  \Delta T = \qty{8}{\celsius} & Temperature drop across specimen
\end{givens}

\begin{assumptions}
  \item Perfectly insulated outer edges
\end{assumptions}

\begin{derivation}
  As the heater is always heating two samples at a time the power of the heater is split evenly between the two samples:
  \begin{equation*}
    \dot{Q}_{\text{sample}} = \frac{\dot{Q}_{\text{heater}}}{2} = \qty{12,5}{\W} 
  .\end{equation*}
  This flow of $\dot{Q}_{\text{sample}} = \qty{12,5}{\W}$ is what sustains the temperature of the samples.
  The thermal conductivity of a material is defined from:
  \begin{equation*}
    \dot{Q}_{\text{cond}} = - kA \frac{\Delta T}{\Delta x}
  .\end{equation*}
  Hence,
  \begin{equation*}
    k = \frac{\dot{Q}_{\text{cond}} \cdot h}{A \cdot \Delta T} = \frac{\qty{12,5}{\W} \cdot \qty{0,5}{\cm}}{\qty{10}{\cm} \cdot \qty{10}{\cm} \cdot \qty{8}{\K} }  = \qty{0,78}{\W\per\m\per\K} 
  .\end{equation*}
\end{derivation}

\begin{finalresult}
  Hence the thermal conductivity of each sample is:
  \begin{equation*}
    k = \qty{0,78}{\W\per\m\per\K} 
  .\end{equation*}
\end{finalresult}



\begin{problemwithimage}{p1-100}{1--100}
  The inner and outer surfaces of a \qty{25}{\cm} thick wall in summer are at \qty{27}{\celsius}  and \qty{44}{\celsius}, respectively.
  The outer surface of the wall exchanges heat by radiation with surrounding surfaces at \qty{40}{\celsius}, and convection heat transfer coefficient of \qty{8}{\W\per\m\squared\per\K}.
  Solar radiation is incident on the surface at a rate of \qty{150}{\W\per\m\squared}.
  If both the emissivity and the solar absorptivity of the outer surface are \num{0,8}, determine the effective thermal conductivity of the wall.
\end{problemwithimage}

\begin{givens}
  w = \qty{25}{\cm}  & Thickness of wall \\
  T_{\text{inner}} = \qty{27}{\celsius} & Inner temperature of wall \\
  T_{\text{outer}} = \qty{44}{\celsius} & Outer temperature of wall \\
  T_{\infty, \text{outer}} = \qty{40}{\celsius} & Ambient temperature outside wall \\
  h = \qty{8}{\W\per\m^2\per\K} & Heat transfer coefficient \\
  \dot{q}_{\text{indicent}} = \qty{150}{\W\per\m\squared} & Solar radiation incident on outer wall \\
  \epsilon = \alpha = \num{0,8} & Solar emissivity and absorptivity of wall
\end{givens}

\begin{assumptions}
  \item The system has reached steady-state
\end{assumptions}

\begin{derivation}
  The radiation heat transfer rate on a solid surface is:
  \begin{equation*}
    \dot{q}_{\text{rad}} = \dot{q}_{\text{emit}} - \dot{q}_{\text{abs}}
  .\end{equation*}
  Here, the net emitted radiation can be found using Boltzmann's law:
  \begin{equation*}
    \dot{q}_{\text{emit}} = \epsilon \sigma \left( T_{s}^{4} - T_{\infty}^{4} \right) = \num{0,8} \cdot \qty{5,670e-8}{\W\per\m\squared\per\K\tothe{4}} \cdot \left( \left( \qty{44}{\celsius} \right)^4 - \left( \qty{40}{\celsius}  \right)^4 \right) = \qty{22,7}{\W\per\m\squared} 
  .\end{equation*}
  The absorbed radiation is found as:
  \begin{equation*}
    \dot{q}_{\text{abs}} = \alpha \cdot \dot{q}_{\text{incident}} = \num{0,8} \cdot \qty{150}{\W\per\m\squared}  = \qty{120}{\W\per\m\squared} 
  .\end{equation*}
  Hence, the radiation heat transfer rate on the outside of the wall is:
  \begin{equation*}
    \dot{q}_{\text{rad}} = \qty{22,7}{\W\per\m\squared} - \qty{120}{\W\per\m\squared} = \qty{97,3}{\W\per\m\squared} 
  .\end{equation*}
  
  From Newton's law of cooling, the convection heat transfer rate is:
  \begin{equation*}
    \dot{q}_{\text{conv}} = h \left( T_{\text{outer}} - T_{\infty} \right) = \qty{8}{\W\per\m\squared\per\K}  \cdot \left( \qty{44}{\K} - \qty{40}{\K}  \right)  = \qty{31}{\W\per\m\squared} 
  .\end{equation*}
  
  Hence, the outside of the wall is gaining heat at a rate of:
  \begin{equation*}
    \dot{q}_{\text{outer}} = \dot{q}_{\text{rad}} - \dot{q}_{\text{conv}} = \qty{97,3}{\W\per\m\squared} - \qty{31}{\W\per\m\squared}  = \qty{66,3}{\W\per\m\squared} 
  .\end{equation*}
  
  Now, the thermal conductivity is easily found as:
  \begin{equation*}
    k = \frac{\dot{q}_{\text{outer}} \cdot w}{ T_{\text{outer}} - T_{\text{inner}}} = \frac{\qty{66,3}{\W\per\m\squared} \cdot \qty{25}{\cm} }{\qty{44}{\celsius} -\qty{27}{\celsius} } = \qty{0,975}{\W\per\m\per\K}
  .\end{equation*}
\end{derivation}

\begin{finalresult}
  Hence, the thermal conductivity of the wall is
  \begin{equation*}
    k = \qty{0,975}{\W\per\m\per\K} 
  .\end{equation*}
  
\end{finalresult}

\begin{problemwithimage}{p1-106}{1--106}
  A flat-plate solar collector is used to heat water by having water flow through tubes attached at the back of the thin solar absorber plate.
  The absorber plate has a surface area of \qty{2}{\m\squared} with emissivity and absorptivity of \num{0,9}.
  The surface temperature of the absorber is \qty{35}{\celsius}, and solar radiation is incident on the absorber at \qty{500}{\W\per\m\squared} with a surrounding temperature of \qty{0}{\celsius}.
  Convection heat transfer coefficient at the absorber surface is \qty{5}{\W\per\m\squared\per\K}, while the ambient temperature is \qty{25}{\celsius}.
  Net heat rate absorbed by the solar collector heats the water from an inlet temperature ($T_{\mathrm{in}}$) to an outlet temperature ($T_{\mathrm{out}}$).
  If the water flow rate is \qty{5}{\g\per\s} with specific heat of \qty{4,2}{\kJ\per\kg\per\K}, determine the temperature rise of the water.
\end{problemwithimage}

\begin{assumptions}
  \item The system has reached steady-state
\end{assumptions}

\begin{givens}
  A_s = \qty{2}{\m\squared}  & Surface area of absorber \\
  \epsilon = \alpha = \num{0,9} & Emissivity and absorptivity of absorber \\
  T_{\text{absorber}} = \qty{35}{\celsius} & Temperature of absorber \\
  \dot{q}_{\text{incident}} = \qty{500}{\W\per\m\squared} & Solar radiation incident at absorber \\
  T_{\text{rad}, \infty} = \qty{0}{\celsius} & Effective surrounding temperature for radiation \\
  h = \qty{5}{\W\per\m\squared\per\K} & Heat transfer coefficient \\
  T_{\text{conv.}, \infty} = \qty{25}{\celsius} & Ambient temperature for convection \\
  \dot{m} = \qty{5}{\g\per\s} & Mass flow rate of water \\
  c_p = \qty{4,2}{\kJ\per\kg\per\K} & Specific heat capacity of water.
\end{givens}

\begin{derivation}
  For radiation we use the same procedure as above and get:
  \begin{equation*}
    \dot{q}_{\text{rad}} = \epsilon \sigma \left( T_{\text{absorber}}^{4} - T_{\text{rad}, \infty}^{4} \right) = \num{0,9} \cdot \qty{5,670e-8}{\W\per\m\squared\per\K\tothe{4}} \cdot \left( \left( \qty{35}{\celsius}  \right)^{4} - \left( \qty{0}{\celsius}  \right)^{4} \right) = \qty{176,05}{\W\per\m\squared} 
  .\end{equation*}
  For convection, we also use the same procedure as above, which yields:
  \begin{equation*}
    \dot{q}_{\text{conv}} = h \left( T_{\text{absorber}} - T_{\text{conv.}, \infty} \right) = \qty{5}{\W\per\m\squared\per\K} \cdot \left( \qty{35}{\celsius} - \qty{25}{\celsius}  \right) = \qty{50}{\W\per\m\squared} 
  .\end{equation*}
  Hence, the total energy balance is:
  \begin{equation*}
    \dot{q}_{\text{absorber}} = \dot{q}_{\text{rad}} - \dot{q}_{\text{conv}} = \qty{176,05}{\W\per\m\squared}  - \qty{50}{\W\per\m\squared}  = \qty{126,05}{\W\per\m\squared} 
  .\end{equation*}
  This means the total absorbed power over the entire area of the absorber is:
  \begin{equation*}
    \dot{Q}_{\text{absorber}} = \dot{q}_{\text{absorber}} \cdot A_s = \qty{126,05}{\W\per\m\squared} \cdot \qty{2}{\m\squared}  = \qty{252.1}{\W}
  .\end{equation*}
  
  As the system is assumed to be at steady state none of this heat will go to warming the absorber -- instead all the $\dot{Q}_{\text{absorber}} = \qty{252,1}{\W}$ goes towards heating the water.
  Using the formula for temperature change as a result of added thermal energy we arrive at:
  \begin{align*}
    \dot{Q}_{\text{absorber}} &= \dot{m} c_p \cdot \Delta T \\
    \Delta T &= \frac{\dot{Q}_{\text{absorber}}}{\dot{m} \cdot c_p} \\
    &= \frac{\qty{252,1}{\W}}{\qty{5}{\g\per\s} \cdot \qty{4,2}{\kJ\per\kg\per\K}} \\
    &= \qty{12}{\K}
  .\end{align*}
\end{derivation}

\begin{finalresult}
  Hence the temperature of the water will increase by:
  \begin{equation*}
    \Delta T = \qty{12}{\K}
  \end{equation*}
  when passing through the absorber.
\end{finalresult}
